\newpage
\begin{center}
  \textbf{\large АННОТАЦИЯ}
\end{center}


\setcounter{page}{2}
Дипломная работа посвящена разработке и применению методики анализа защищённости мобильных приложений с использованием прокси-сервера для перехвата и анализа сетевых вызовов. В условиях стремительного роста популярности мобильных приложений и увеличения объёма передаваемых данных вопросы обеспечения их безопасности становятся особенно актуальными. Целью работы является исследование уязвимостей мобильных приложений путём анализа их сетевого взаимодействия и разработка рекомендаций по повышению уровня их защищённости, а также использование полученных данных для создания отечественного аналога приложения.

В рамках исследования была разработана методика, основанная на использовании прокси-сервера для перехвата и анализа сетевых запросов между мобильным приложением и сервером. Данный подход был протестирован на трёх различных приложениях: приложение для заказа еды, приложение для доступа к бизнес-клубу с годовым оборотом от 100 миллионов рублей и игра Minecraft для сбора данных об игроках на общем сервере. В ходе анализа были выявлены особенности сетевого взаимодействия, включая типы передаваемых данных, протоколы и потенциальные уязвимости, связанные с недостаточной защитой соединений или некорректной обработкой данных.

Полученные результаты позволили не только выявить уязвимости в исследуемых приложениях, но и использовать собранные данные для разработки отечественного аналога приложения для заказа еды. В работе подробно описаны этапы анализа, включая настройку прокси-сервера, сбор и обработку данных, а также выявление характерных уязвимостей, таких как отсутствие шифрования, слабая аутентификация или утечка конфиденциальной информации. Кроме того, на основе собранных данных были предложены рекомендации по повышению защищённости мобильных приложений, включая внедрение современных протоколов шифрования, усиление механизмов аутентификации и минимизацию объёма передаваемых данных.

Разработанная методика может быть применена для анализа защищённости других мобильных приложений, а также для создания более безопасных программных продуктов. Результаты работы представляют практическую ценность для разработчиков мобильных приложений, специалистов по информационной безопасности и исследователей, занимающихся анализом сетевых взаимодействий. Работа способствует развитию отечественных технологий в области мобильной разработки и повышению уровня кибербезопасности.


\newpage
\renewcommand{\contentsname}{\centerline{\large СОДЕРЖАНИЕ}}
\tableofcontents

\newpage
\begin{center}
  \textbf{\large ВВЕДЕНИЕ}
\end{center}
\addcontentsline{toc}{chapter}{ВВЕДЕНИЕ}


\textbf{Актуальность}

В современном мире мобильные приложения стали неотъемлемой частью повседневной жизни, обеспечивая доступ к широкому спектру услуг — от заказа еды и управления финансами до социальных взаимодействий и развлечений. С ростом популярности мобильных технологий стремительно увеличивается объем передаваемой информации, включая персональные данные пользователей, финансовые транзакции и конфиденциальную корпоративную информацию. Это делает мобильные приложения привлекательной целью для кибератак, таких как перехват данных, несанкционированный доступ и эксплуатация уязвимостей.

Актуальность проблемы анализа защищенности мобильных приложений обусловлена несколькими факторами. Во-первых, недостаточная защита данных в мобильных приложениях может привести к серьезным последствиям, включая утечку конфиденциальной информации, финансовые потери и репутационный ущерб для компаний. Во-вторых, с развитием технологий, таких как прокси-серверы, злоумышленники получают новые возможности для анализа и эксплуатации сетевых взаимодействий приложений. В-третьих, в условиях цифровизации и импортозамещения в России разработка отечественных аналогов популярных приложений требует особого внимания к вопросам безопасности, чтобы обеспечить конкурентоспособность и доверие пользователей.

Настоящая работа посвящена исследованию методик анализа защищенности мобильных приложений, включая использование прокси-серверов для перехвата и анализа сетевых вызовов. В рамках исследования были протестированы приложения различного назначения — от сервисов заказа еды до игровых платформ, таких как Minecraft, — что позволило выявить уязвимости и собрать данные для создания защищенных отечественных аналогов. Изучение данных методик способствует повышению уровня безопасности мобильных приложений и разработке эффективных решений для защиты пользовательских данных.

\newpage

\textbf{Цель работы} -- разработка универсальной методики анализа защищенности мобильных приложений, обеспечивающей выявление уязвимостей, проверку воспроизводимости на различных типах приложений и разработку защитных мер, а также создание отечественных аналогов приложений с закрытым исходным кодом.

\textbf{Задачи работы:}
\begin{enumerate}
\item Провести анализ уязвимостей мобильных приложений с клиент-серверной архитектурой, включая изучение механизмов обмена данными и выявление рисков, связанных с отсутствием асимметричного шифрования. 
\item Разработать и проверить универсальную методику проведения атак типа "человек посередине" на общедоступных приложениях.
\item Сформировать рекомендации по предотвращению атак и безопасному обращению с данными для повышения защищенности приложений.
\item Исследовать возможность разработки отечественных аналогов приложений с открытым исходным кодом на основе анализа коммуникации закрытых приложений.
\end{enumerate}


\textbf{Научной новизной обладают следующие результаты работы:}
\begin{enumerate}
\item Установлено, что отсутствие асимметричного шифрования и доверительных сертификатов в большинстве гражданских мобильных приложений с клиент-серверной архитектурой позволяет проводить тривиальные атаки типа "человек посередине", создавая значительные риски для безопасности данных.
\item Выявлена универсальность разработанной методики анализа уязвимостей, подтверждённая её успешным применением к различным типам приложений, включая системы удалённого заказа еды, приложения бизнес-клубов и низкоуровневые TCP-протоколы.
\item Определены ключевые недостатки современных коммерческих приложений в области безопасного обмена данными, что подчеркивает необходимость внедрения стандартов безопасности, таких как OWASP, в процесс разработки.
\end{enumerate}


\textbf{Апробация} основных результатов работы проводилась на следующих конференциях:
\begin{enumerate}
\item Underconf конференция сообщества информационной безопасности 29 сентября 2024г.
\item IT Purple Conf 15 марта 2025г.
\item 67-я Всероссийская научная конференция МФТИ 3 апреля 2025г.
\end{enumerate}

